\chapter{An explicit ancient oval}\label{sec:examples-ancient-curves}
An ancient solution exists for every time before some finite time. It need not preserve its shape. The explicit oval below is often called the Angenent oval or paperclip solution; \cite{BV11}, Section 3.1, equation (14), records the formula and discusses its provenance. We verify the evolution and its limiting models directly. The earlier shrinking circle and grim reaper appear here as two different limits of one solution.

\begin{theorem}[The implicit oval is an ancient curve-shortening flow]\label{thm:explicit-ancient-oval}
For $t<0$, the set
\[
 \Gamma_t=\{(x,y):|x|<\pi/2,\ e^t\cosh y=\cos x\}
\]
is a smooth embedded strictly convex closed curve. These curves solve curve shortening flow and shrink to the origin as $t\uparrow0$. Their coordinate extrema are
\[
 |x|\leq\arccos(e^t),\qquad |y|\leq\operatorname{arcosh}(e^{-t}).
\]
\end{theorem}
\begin{proof}
Set $F(x,y,t)=e^t\cosh y-\cos x$. On the strip $|x|<\pi/2$, its spatial Hessian is positive definite. Its negative sublevel is nonempty since $F(0,0,t)<0$, and the displayed bounds show its closure stays in a compact subset of the strip. It is a strictly convex bounded domain. On its boundary,
\[
 |\nabla F|^2=\sin^2x+e^{2t}\sinh^2y=1-e^{2t}>0,
 \qquad D^2F=\cos x\,I.
\]
Thus the boundary is smooth and has outward curvature $\kappa=\operatorname{div}(\nabla F/|\nabla F|)=\cos x/\sqrt{1-e^{2t}}>0$. Differentiating $F=0$ in time gives outward normal velocity $-F_t/|\nabla F|=-\cos x/\sqrt{1-e^{2t}}=-\kappa$. This is curve shortening. The coordinate bounds follow by setting the other coordinate to zero and tend to zero as $t\uparrow0$. The gradient computation avoids treating the two graph branches as separate curves at their endpoints. Compare \cite{BV11}, equation (14), p.~7.
\end{proof}

\begin{proposition}[The ancient oval changes shape]\label{prop:ancient-oval-curvature-ratio}
The extreme curvatures of $\Gamma_t$ satisfy
\[
 \kappa_{\max}(t)=\frac1{\sqrt{1-e^{2t}}},\qquad
 \kappa_{\min}(t)=\frac{e^t}{\sqrt{1-e^{2t}}},\qquad
 \frac{\kappa_{\max}}{\kappa_{\min}}=e^{-t}.
\]
In particular this ancient solution is not a soliton moving by Euclidean similarities. Its curvature becomes asymptotically uniform as extinction approaches, but is very nonuniform far backward in time.
\end{proposition}
\begin{proof}
The curvature formula from Theorem \ref{thm:explicit-ancient-oval} is proportional to $\cos x$. On the curve this ranges from $e^t$ at the horizontal extrema to one at the vertical extrema. This proves all formulas. A Euclidean similarity changes all curvatures by the same positive scale factor, and so preserves their maximum-to-minimum ratio. Here that ratio changes strictly with time. Its limits are one as $t\uparrow0$ and infinity as $t\to-\infty$.
\end{proof}

\begin{proposition}[Area and the shrinking-circle limit]\label{prop:ancient-oval-extinction-model}
The area enclosed by $\Gamma_t$ is $-2\pi t$. The rescaled curves $(-t)^{-1/2}\Gamma_t$ converge smoothly as $t\uparrow0$ to the circle of radius $\sqrt2$ centered at zero.
\end{proposition}
\begin{proof}
The normal velocity is $-\kappa$, so the area derivative is $-\int_{\Gamma_t}\kappa\,ds=-2\pi$, using the turning-angle formula for a smooth simple convex curve. The enclosed area tends to zero by the coordinate bounds in Theorem \ref{thm:explicit-ancient-oval}; integration gives $-2\pi t$. This is the area calculation in \cite{WhiMCF15}, Theorem 1.1, applied to the explicit oval.

For smooth convergence put $s=-t>0$ and $(x,y)=\sqrt s(X,Y)$. On bounded sets the rescaled defining function satisfies
\[
 \frac{e^{-s}\cosh(\sqrt sY)-\cos(\sqrt sX)}s
 =\frac{X^2+Y^2}{2}-1+O(s)
\]
with convergence of every spatial derivative, by the power series. The rescaled coordinate extrema tend to $\sqrt2$, so all curves stay in a fixed bounded set. Strict convexity and inclusion of the origin inside each oval give a single radial intersection in every direction. The limiting function has nonzero radial derivative on the radius-$\sqrt2$ circle. The implicit function theorem, uniformly over that compact circle, gives smooth convergence of the radial functions and hence of the curves. The normalization agrees with the shrinker radius in Example \ref{ex:soliton-model-comparison}.
\end{proof}

\begin{proposition}[The two grim-reaper limits backward in time]\label{prop:ancient-oval-translator-limits}
On every compact subinterval of $(-\pi/2,\pi/2)$, the upper and lower graph branches of the ancient oval satisfy, as $t\to-\infty$,
\[
 y_+(x,t)+t-\log2\longrightarrow\log\cos x,
 \qquad y_-(x,t)-t+\log2\longrightarrow-\log\cos x
\]
smoothly in $x$. These are the downward and upward grim-reaper profiles, respectively. The convergence is local in the horizontal strip, not uniform through its limiting endpoints.
\end{proposition}
\begin{proof}
The branches are $y_\pm=\pm\operatorname{arcosh}(e^{-t}\cos x)$. For large $z$, the identity $\operatorname{arcosh}z=\log(z+\sqrt{z^2-1})$ gives $\operatorname{arcosh}z=\log(2z)+O(z^{-2})$. On a compact subinterval the factor $\cos x$ stays positive, so substitution gives an $O(e^{2t})$ error with every fixed number of $x$ derivatives. This proves both limits. The signs of their vertical motions agree with Proposition \ref{prop:grim-reaper-speed-family} and reflection across the horizontal axis. This is the backward translator behavior discussed in \cite{BV11}, p.~7; the error bound here is derived from the exact inverse-hyperbolic-cosine expression.
\end{proof}
